
\begin{lemme}
\label{XIII.4.2}
Sean $S$ un esquema conexo, $f\colon X\to S$ un morfismo localmente
$0$-acíclico y $0$-acíclico, y~$\LL$ un conjunto de números primos. Supongamos
\marginpar{417}
que, para todo haz constante finito de~$\LL$-grupos~$F$ sobre~$X$,
$(F,f)$ es cohomológicamente propio relativamente a~$S$ en dimensión
$\leqslant1$, y que, para todo recubrimiento étale~$S'$ de~$S$ y para todo
recubrimiento étale~$E$ de~$X'$, cociente de un recubrimiento galoisiano
cuyo grupo es un~$\LL$-grupo, $f'_*E$ es constructible. Entonces, si~$\bar{s}$
es un punto geométrico de~$S$ y~$a$ un punto geométrico de la fibra
$X_{\bar{s}}$, la sucesión de homomorfismos de grupos
$$
\pi_1^{\LL}(X_{\bar{s}},a)\to\pi_1'(X,a)\to\pi_1(S,a)\to1
$$
es exacta.
\end{lemme}

Se satisfacen las hipótesis de~\Ref{XIII.4.1}. En efecto, de
\Ref{XIII.1.13}~3) se sigue que, para todo esquema~$S'$ étale sobre~$S$ y
para todo recubrimiento étale~$E$ de~$X'$, cociente de un recubrimiento
galoisiano cuyo grupo es un~$\LL$-grupo, $(E,f')$ es cohomológicamente propio
relativamente a~$S'$ en dimensión $\leqslant 0$.

\begin{proposition}
\label{XIII.4.3}
\setcounter{subsection}{\value{proposition}}% sincroniza la etiqueta visible XIII.4.3.1
Sean $S$ un esquema conexo, $\LL$ un conjunto de números primos,
$f\colon X\to S$ un morfismo $0$-acíclico y localmente $1$-asférico
para~$\LL$ \textup{(SGA~4 XV~1.11)}, y~$g\colon S\to X$ una sección de~$f$.
Sean~$\bar{s}$ un punto geométrico de~$S$ y~$a$ un punto geométrico de la
fibra~$X_{\bar{s}}$. Supongamos que, para todo haz de
\ifthenelse{\boolean{orig}}
{$\LL$-grupo}
{$\LL$-grupos}
constante~$F$, $(F,f)$ es cohomológicamente propio en dimensión
$\leqslant1$, que la imagen directa por~$f$ del stack de los torsores bajo
$F$ es un stack $1$-constructible \eqref{XIII.0} y que,
\ifthenelse{\boolean{orig}}
{para todo recubrimiento étale~$E$ de~$X'$}
{para todo recubrimiento étale~$S'$ de~$S$ y para todo
recubrimiento étale~$E$ de~$X'$}, cociente de un recubrimiento galoisiano
cuyo grupo es un~$\LL$-grupo, $f'_*E$ es constructible. Entonces, la sucesión
de homomorfismos de grupos
\begin{equation*}
\label{eq:XIII.4.3.1}
\tag{\thesubsection.1}
1\to\pi_1^{\LL}(X_{\bar{s}},a)\lto{u} \pi_1'(X,a)\lto{v}\pi_1(S,a)\to1
\end{equation*}
es exacta.
\end{proposition}

Habida cuenta de~\Ref{XIII.4.2}, basta demostrar la inyectividad del
morfismo~$u$, es decir, probar que, para todo recubrimiento principal
$\overline{Z}$ de~$X_{\bar{s}}$
\marginpar{418}
cuyo grupo es un~$\LL$-grupo~$C$, existen un recubrimiento étale~$Z$ de~$X$
y un morfismo de una componente conexa de~$Z|X_{\bar{s}}$ hacia
$\overline{Z}$ (\Ref{V.6.8}). Sean, pues, $\overline{Z}$ un recubrimiento
principal de~$X_{\bar{s}}$ cuyo grupo es un~$\LL$-grupo~$C$ y~$\bar{z}$ su
clase en~$\H^1(X_{\bar{s}},C_{X_{\bar{s}}})$. Según \Ref{XIII.1.5}~d), se
tiene un isomorfismo canónico
$$
\left(\R^1f_*C_X\right)_{\bar{s}}\isomto\H^1(X_{\bar{s}},C_{X_{\bar{s}}}),
$$
y, según~\Ref{XIII.1.16}, $\R^1f_*C_X$ es un haz localmente constante y
constructible. Puede encontrarse, por tanto, un recubrimiento étale~$S'$ de~$S$
tal que~$\R^1f_*C_X|S'$ sea constante. Si~$\bar{s}\to S'$ es un punto
geométrico situado sobre el punto geométrico~$\bar{s}\to S$, existe un
elemento~$z$ de~$\H^0(S', \R^1f_*C_X)$ cuya imagen en
$\H^1(X_{\bar{s}},C_{X_{\bar{s}}})$ es~$\bar{z}$. Según el lema
\Ref{XIII.4.3.1} que sigue, puede encontrarse un recubrimiento étale~$S'_1$
de~$S'$ y un torsor~$P$ sobre~$X'_1$ bajo~$C$ cuya imagen en
$\H^0(S'_1,\R^1f_*C)$ sea igual a la restricción de~$z$. El torsor~$P$ es
representable por un recubrimiento étale~$Z$ de~$X'_1=X\times_SS'_1$ tal que
$Z\times_{X'_1}X_{\bar{s}}$ sea isomorfo a~$\overline{Z}$. Si se considera
$Z$ como un recubrimiento étale de~$X$, se tiene entonces un morfismo
de~$Z\times_XX_{\bar{s}}$ hacia~$\overline{Z}$, lo que concluye la
demostración.
